\notepage{47}{Integer and half-integer angular momentum}
The endpoint result gives
\[
 \alpha=\ell\hbar(\ell\hbar+\hbar)=\ell(\ell+1)\hbar^2,
 \qquad \beta=m\hbar,
\]
where $m=-\ell,-\ell+1,\ldots,\ell$ in steps of one. Relabeling the states,
\[
 L^2\ket{\ell,m}=\ell(\ell+1)\hbar^2\ket{\ell,m},\qquad
 L_z\ket{\ell,m}=m\hbar\ket{\ell,m}.
\]
The algebra allows $\ell=0,\tfrac12,1,\tfrac32,\ldots$.

\smallheading{The orbital coordinate representation}
Starting instead from $L_z=-\ii\hbar\partial_\phi$,
\[
 -\ii\hbar\frac{\partial}{\partial\phi}
 \braket{\phi}{\alpha,\beta}
 =\beta\braket{\phi}{\alpha,\beta},
 \qquad
 \braket{\phi}{\alpha,\beta}\propto\ee^{\ii\beta\phi/\hbar}.
\]
For a scalar orbital wavefunction, single-valuedness under
$\phi\mapsto\phi+2\pi$ requires
\[
 \ee^{2\pi\ii\beta/\hbar}=1\quad\Longrightarrow\quad
 \beta=m\hbar,\qquad m\in\mathbb Z.
\]
Thus the orbital representation restricts $\ell$ and $m$ to integers,
whereas the abstract commutation algebra also admits half-integers.

\smallheading{Why the algebra is more general}
The total angular momentum is
\[
 \mathbf J=\mathbf L\otimes I_s+I_r\otimes\mathbf S,
\]
usually abbreviated as $\mathbf J=\mathbf L+\mathbf S$.
Spin satisfies $\mathbf S\times\mathbf S=\ii\hbar\mathbf S$;
because spin and orbital operators act on different factors,
\[
 \mathbf J\times\mathbf J=\ii\hbar\mathbf J.
\]
Consequently, deriving eigenvalues only from this algebra does not
specify that the operator must be $\mathbf R\times\mathbf P$.
It includes spin and total angular momentum, for which half-integer
representations occur.
